% SOURCE_AUTHORITY: sga5_fr_workpass.tex, lines 13205--13229 (LF-normalized unit).
% SOURCE_SHA256: 791F4EFFC5E02832D5D77ED03518C8156D6F07E4C8238B03545DB93D883FBB28
Por último, para establecer (7.13) basta utilizar la fórmula de Euler--Poincaré (7.2) para el complejo de haces $i_!(F)$:
\[
\EP_\Delta(V,F)
=
\clK{\Delta}(\RG_{!V}(F))
=
\clK{\Delta}(\RG_C(i_!(F)))
=
\EP(C)\clK{\Delta}(F_{\bar\eta})
-\sum_{x\in C(k)}\varepsilon_x^\Delta(i_!(F)).
\]
Ahora bien,
\[
\sum_{x\in C(k)}\varepsilon_x^\Delta(i_!(F))
=
\sum_{x\in V(k)}\varepsilon_x^\Delta(F)
+\sum_{x\in Z}\varepsilon_x^\Delta(F).
\]
Pero para $x\in Z$ se tiene $(i_!(F))_{\bar x}=0$ y, por consiguiente, para tal $x$ se tiene
\[
\varepsilon_x^\Delta(i_!(F))
=
\alpha_x^\Delta(F)+\clK{\Delta}(F_{\bar\eta}),
\]
de donde resulta la fórmula (7.13).
\end{prooftext}
